\documentclass[11pt,a4paper]{article}

% --- 核心宏包与中文支持 ---
\usepackage[utf8]{inputenc}
\usepackage{xeCJK}        % 必须使用 XeLaTeX 编译器
\usepackage{amsmath, amssymb, amsfonts}
\usepackage{listings}
\usepackage{xcolor}
\usepackage{geometry}
\usepackage{tcolorbox}
\usepackage{booktabs}
\usepackage{tabularx}
\usepackage{setspace}
\usepackage{enumitem}
\usepackage[hidelinks]{hyperref}
\usepackage{bookmark}
\usepackage{fancyhdr}

% --- PDF 书签与元数据 ---
\hypersetup{
    colorlinks=false,
    pdfborder={0 0 0},
    bookmarks=true,
    bookmarksnumbered=true,
    bookmarksopen=true,
    bookmarksopenlevel=2,
    pdfcreator={XeLaTeX},
    pdftitle={04 - Modules and Packages},
    pdfauthor={黄子扬},
}

% --- 页面美化设置 ---
\geometry{top=0.7in,bottom=0.8in,left=0.6in,right=0.6in}
\onehalfspacing
\tcbuselibrary{skins,breakable}
\setlength{\parindent}{0pt}
\setlength{\parskip}{0.6em}

% --- 代码样式配置 ---
\definecolor{mainblue}{RGB}{0,51,102}
\definecolor{examred}{RGB}{180,0,0}
\definecolor{codebg}{RGB}{248,248,248}
\definecolor{commentgreen}{RGB}{0,128,0}

\lstset{
    backgroundcolor=\color{codebg},
    basicstyle=\ttfamily\footnotesize,
    keywordstyle=\color{blue}\bfseries,
    commentstyle=\color{commentgreen},
    stringstyle=\color{orange},
    breaklines=true,
    showstringspaces=false,
    frame=leftline,
    xleftmargin=2em,
    numbers=left,
    numbersep=5pt,
    numberstyle=\tiny\color{gray},
    language=Python,
    morekeywords={None, True, False, import, as, from, np, pd, def, return, for, in, range}
}

% --- 自定义教学框 ---
\newtcolorbox{concept}[1]{colback=blue!5, colframe=mainblue, fonttitle=\bfseries, title={{【Concept / 核心概念】#1}}, arc=3pt, breakable}
\newtcolorbox{warning}[1]{colback=red!5, colframe=examred, fonttitle=\bfseries, title={{【Warning / 避坑指南】#1}}, arc=3pt, breakable}
\newtcolorbox{examplebox}[1]{colback=green!2, colframe=green!40!black, fonttitle=\bfseries, title={{【Example / 实操案例】#1}}, arc=2pt, breakable}

% --- 页眉页脚 ---
\pagestyle{fancy}
\fancyhf{}
\fancyhead[L]{\small\bfseries 04 - Modules and Packages / 模块与包}
\fancyhead[R]{\small\thepage}
\renewcommand{\headrulewidth}{0.5pt}

% --- 文档信息 ---
\title{\Huge \bfseries 04 - Modules and Packages}
\author{\Large \textbf{哈尔滨工业大学 \quad 黄子扬}}
\date{2026年6月8日}

\begin{document}

\maketitle
\tableofcontents
\newpage

% ======================================================================
\section{Modules and Importing / 模块与导入}
% ======================================================================

\begin{concept}{What is a Module? / 什么是模块？}
    A module is a file containing Python code that can be re-used in other Python code files. \\
    模块是一个包含 Python 代码的文件，可以在其他代码中重复使用。

    There are libraries already included with Python and others we have to install before being able to import and use them.
\end{concept}

\subsection{Four Import Variations / 四种导入方式}

\begin{concept}{Import Syntax / 导入语法}
    Check the following four variations for importing modules from packages:
    \begin{enumerate}[itemsep=0.5em]
        \item \texttt{import [package]} — imports the entire package.
        \item \texttt{import [package] as [other\_name]} — imports with an alias (shorter name).
        \item \texttt{from [package] import [module]} — imports only a specific module/function.
        \item \texttt{from [package] import [module] as [other\_name]} — imports a specific item with alias.
    \end{enumerate}
\end{concept}

\begin{examplebox}{Import Examples / 导入示例}
\begin{lstlisting}
# 1. Basic import
import numpy
arr = numpy.array([1, 2, 3])

# 2. Import with alias (most common for NumPy/Pandas)
import numpy as np
arr = np.array([1, 2, 3])

# 3. Import specific function
from math import sqrt
print(sqrt(16))             # 4.0

# 4. Import specific function with alias
from math import sqrt as square_root
print(square_root(16))      # 4.0
\end{lstlisting}
\end{examplebox}

% ======================================================================
\section{Standard Library / 标准库核心模块}
% ======================================================================

\begin{concept}{What is the Standard Library? / 什么是标准库？}
    The \textbf{Standard Library} doesn't need to be installed separately — it comes with your Python installation. These are often called \textbf{inbuilt packages}. \\
    标准库随 Python 安装自带，无需额外安装。也叫"内置包"。

    External libraries like NumPy and Pandas must be installed before using them.
\end{concept}

\subsection{Copy Module / 复制模块}

\begin{concept}{Assignment vs Copy / 赋值与复制的区别}
    Assignment (\texttt{=}) does not copy objects; it creates a \textbf{binding} (reference) between a name and an object. For mutable collections, a copy is sometimes needed so one can change one copy without changing the other. \\
    赋值只是创建引用（别名），不是真正的复制。对于可变集合，有时需要真正的副本来独立修改。

    \begin{itemize}[itemsep=0.5em]
        \item \textbf{Shallow Copy / 浅拷贝}（\texttt{copy.copy}）：Creates a new container, but elements inside still reference the original objects.
        \item \textbf{Deep Copy / 深拷贝}（\texttt{copy.deepcopy}）：Recursively copies everything — completely independent from the original.
    \end{itemize}
\end{concept}

\begin{examplebox}{Shallow Copy Example / 浅拷贝案例}
\begin{lstlisting}
import copy
animals = ["lion", "tiger"]

# Without copy — just a reference
large_cats = animals
large_cats.append("Leopard")
print(animals)      # ['lion', 'tiger', 'Leopard'] — CHANGED!

# With copy.copy — independent copy
animals = ["lion", "tiger"]
large_cats = copy.copy(animals)
large_cats.append("Leopard")
print(large_cats)   # ['lion', 'tiger', 'Leopard']
print(animals)      # ['lion', 'tiger'] — unchanged!
\end{lstlisting}
\end{examplebox}

\subsection{Math Module / 数学模块}

\begin{concept}{What's Inside math? / math 模块的内容}
    The \texttt{math} package includes mathematical constants and functions:
    \begin{itemize}[itemsep=0.3em]
        \item \textbf{Constants / 常量}：\texttt{math.pi}, \texttt{math.e}.
        \item \textbf{Functions / 函数}：\texttt{log}, \texttt{cos}, \texttt{sin}, \texttt{ceil} (向上取整), \texttt{floor} (向下取整).
    \end{itemize}
\end{concept}

\begin{examplebox}{math Module Examples / math 模块示例}
\begin{lstlisting}
import math

print(math.log(5, 2))       # Logarithm base 2 of 5
print(math.pi)              # 3.141592653589793
print(math.floor(2.6))      # 2 (round down)
print(math.ceil(2.6))       # 3 (round up)
\end{lstlisting}
\end{examplebox}

\subsection{Time Module / 时间模块}

\begin{concept}{What's Inside time? / time 模块的内容}
    The \texttt{time} package deals with time and timing-related things. \\
    \texttt{time.time()} returns the seconds that have passed since \textbf{January 1st 1970 00:00:00} (the Unix epoch).
\end{concept}

\begin{examplebox}{Time Module Examples / time 模块示例}
\begin{lstlisting}
import time

# Get current Unix timestamp / 获取当前时间戳
print(time.time())          # e.g., 1717718400.123456

# Pause execution for 2 seconds / 暂停 2 秒
print("Wait...")
time.sleep(2)
print("Hi")                 # Printed after 2 seconds

# Timing how long a function takes / 测量函数执行时间
def my_function():
    for n in range(1000000):
        n += 2 * n

start_time = time.time()
my_function()
elapsed = time.time() - start_time
print(f"{elapsed:.4f} seconds")
\end{lstlisting}
\end{examplebox}

\begin{examplebox}{Time Object with localtime() / 本地时间对象}
\begin{lstlisting}
now = time.localtime()
print(now)                          # time.struct_time object
print(now.tm_year)                  # Current year, e.g., 2026
print(now.tm_mon)                   # Current month
print(now.tm_mday)                  # Current day
\end{lstlisting}
\end{examplebox}

% ======================================================================
\section{External Libraries \& Pip / 外部库与包管理}
% ======================================================================

\begin{concept}{Why Pip? / 为什么需要 Pip？}
    External packages need to be installed before using them. They are not included with Python by default. \texttt{pip} is the package manager for Python — it automates installing, upgrading, and removing third-party packages. \\
    \texttt{pip} 是 Python 的包管理器，用于安装、升级和卸载第三方包。

    \texttt{pip} is a command line tool — you must run it from a terminal.
\end{concept}

\begin{examplebox}{Common pip Commands / 常用 pip 命令}
\begin{lstlisting}
# Check pip version / 检查 pip 版本
pip --version

# Upgrade pip itself / 升级 pip 自身
pip install --upgrade pip

# List all installed packages / 列出已安装的所有包
pip list

# Install a package / 安装包
pip install numpy

# Install a specific version / 安装特定版本
pip install numpy==1.16.5

# Uninstall a package / 卸载包
pip uninstall numpy
\end{lstlisting}
\end{examplebox}

% ======================================================================
\section{Numerical Computing with NumPy / NumPy 数值计算}
% ======================================================================

\subsection{Why NumPy? / 为什么 NumPy 这么快？}

\begin{concept}{The Type-Checking Bottleneck / 类型检查瓶颈}
    Python lists can store any object, so the interpreter must check the data type for \textbf{every} operation. NumPy uses \textbf{ndarrays} with a \textbf{fixed data type}, skipping these checks and using \textbf{vectorized operations}. \\
    Python 列表因为支持异构数据，每次计算都要检查类型。NumPy 通过固定数据类型（\texttt{dtype}）并使用向量化计算，速度快几十倍。
\end{concept}

\begin{examplebox}{Speed Comparison / 速度对比}
\begin{lstlisting}
import random
import numpy as np

# Create 1,000,000 random integers with Python list
measurements = [random.randint(150, 200) for _ in range(1_000_000)]

# NumPy version
measurements_array = np.array(measurements)

# NumPy's vectorized mean is MUCH faster
np.mean(measurements_array)     # ~50x faster than sum/len on list
\end{lstlisting}
\end{examplebox}

\subsection{Creating NumPy Arrays / 创建 NumPy 数组}

\begin{concept}{Multiple Ways to Create Arrays / 多种创建方式}
    \begin{itemize}[itemsep=0.5em]
        \item \textbf{From list}：\texttt{np.array([10, 2, 35])}
        \item \textbf{arange}：Like \texttt{range()} but returns ndarray. \texttt{np.arange(start=2, stop=14, step=2)}
        \item \textbf{linspace}：Creates array with N evenly spaced values. \texttt{np.linspace(start=-5, stop=5, num=9)}
        \item \textbf{zeros / ones}：Arrays filled with 0 or 1. \texttt{np.zeros(3)}, \texttt{np.ones((2, 3))}
    \end{itemize}
\end{concept}

\begin{examplebox}{Creating Arrays / 创建数组案例}
\begin{lstlisting}
import numpy as np

# From list
print(np.array([10, 2, 35]))            # [10  2 35]

# arange: like range but NumPy
print(list(range(5)))                   # [0, 1, 2, 3, 4]
print(np.arange(start=2, stop=14, step=2))  # [ 2  4  6  8 10 12]

# linspace: N evenly spaced values
print(np.linspace(start=-5, stop=5, num=9))
# [-5.  -3.75 -2.5 -1.25  0.   1.25  2.5  3.75  5. ]

# zeros and ones
print(np.zeros(3))                      # [0. 0. 0.]
print(np.ones(3))                       # [1. 1. 1.]

# Multidimensional with shape argument
print(np.zeros((2, 2)))                 # 2x2 matrix of zeros
# [[0. 0.]
#  [0. 0.]]
print(np.ones((2, 3)))                  # 2x3 matrix of ones
\end{lstlisting}
\end{examplebox}

\subsection{Anatomy of Arrays / 数组解剖}

\begin{concept}{dtype, shape, ndim / 三大核心属性}
    \begin{itemize}[itemsep=0.5em]
        \item \textbf{\texttt{.dtype}}：Data type of the array elements. Can be \texttt{int}, \texttt{float}, \texttt{bool}, etc.
        \item \textbf{\texttt{.shape}}：A tuple with the number of elements in each dimension. E.g., \texttt{(3, 4)} for a 3×4 array.
        \item \textbf{\texttt{.ndim}}：The total number of dimensions.
    \end{itemize}
\end{concept}

\begin{examplebox}{Array Anatomy / 数组属性详解}
\begin{lstlisting}
import numpy as np

# dtype — data type
values = [0, 1, 2, 3, 4]
int_arr = np.array(values, dtype='int')
print(int_arr, int_arr.dtype)           # [0 1 2 3 4] int64

# NumPy casts values to match dtype
bool_arr = np.array(values, dtype='bool')
print(bool_arr, bool_arr.dtype)
# [False  True  True  True  True] bool

# Without explicit dtype: "smallest common denominator"
values = [0, 1, 2.5, 3, 4]
float_arr = np.array(values)
print(float_arr, float_arr.dtype)       # all become float64

# 1D Array
one_dim_arr = np.array([1, 2, 3, 4])
print(one_dim_arr.shape)                # (4,)
print(one_dim_arr.ndim)                 # 1

# 2D Array
values = [[1, 2, 3, 4, 5],
          [1, 2, 3, 4, 1],
          [1, 2, 3, 4, 2]]
two_dim_arr = np.array(values)
print(two_dim_arr.shape)                # (3, 5)
print(two_dim_arr.ndim)                 # 2

# Access 2D elements: [row, column]
print(two_dim_arr[1, 3])                # 4
two_dim_arr[1, 3] = 10                 # Lists are mutable!
print(two_dim_arr)

# 3D Array
values = [[[1, 2, 3, 4]] * 3] * 6
three_dim_arr = np.array(values)
print(three_dim_arr.shape)              # (6, 3, 4)
print(three_dim_arr.ndim)               # 3
print(three_dim_arr[1, 1, 1])           # Access 3D element
\end{lstlisting}
\end{examplebox}

\subsection{Reshape / 重塑数组形状}

\begin{examplebox}{Using reshape() / reshape() 用法}
\begin{lstlisting}
a = np.arange(start=2, stop=14)         # [2, 3, ..., 13] — 12 elements
print(a.shape)                          # (12,)

# Reshape to 3x4
print(a.reshape(3, 4))
# [[ 2  3  4  5]
#  [ 6  7  8  9]
#  [10 11 12 13]]

# -1 auto-detects the size of that dimension
print(a.reshape(2, -1))                 # 2x6
print(a.reshape(2, -1).shape)           # (2, 6)
\end{lstlisting}
\end{examplebox}

\subsection{Comparing Arrays / 数组比较}

\begin{examplebox}{Equality and np.isclose() / 相等比较与 np.isclose()}
\begin{lstlisting}
a = np.zeros((3, 3))
b = np.zeros((3, 3))
print(a == b)                   # All True (element-wise comparison)

# Floating point precision issue!
a[0, 0] += 0.000000000000001
print(a == b)                   # One element is now False!

# Use np.isclose() for floating point comparison
print(np.isclose(a, b))         # All True (tolerates tiny differences)
\end{lstlisting}
\end{examplebox}

\subsection{Mathematical Operations / 数学运算（向量化）}

\begin{concept}{Vectorization / 向量化}
    Numpy's mathematical functions operate on arrays in a \textbf{vectorized} manner — applied to each element without explicit for-loops. Vectorized functions are called \textbf{ufuncs} (universal functions) in NumPy. \\
    向量化操作直接作用于整个数组的每个元素，不需要写 \texttt{for} 循环。
\end{concept}

\begin{examplebox}{Vectorized Operations / 向量化运算}
\begin{lstlisting}
arr = np.arange(1, 10)          # [1 2 3 4 5 6 7 8 9]

# Element-wise operations / 逐元素运算
print(arr * 3)                  # [ 3  6  9 12 15 18 21 24 27]
print(arr * arr)                # [ 1  4  9 16 25 36 49 64 81]
print(arr + (arr * 2))          # [ 3  6  9 12 15 18 21 24 27]
print(arr - arr)                # [0 0 0 0 0 0 0 0 0]
print(arr / arr)                # [1. 1. 1. ...]
print(arr ** 2)                 # [ 1  4  9 16 25 36 49 64 81]

# Matrix multiplication (inner product for 1D)
print(arr @ arr)                # 285 (= 1*1 + 2*2 + ... + 9*9)
# Same as: np.sum(arr * arr)

# Universal functions / 通用数学函数
print(np.log(arr))
print(np.log2(arr))
print(np.exp(arr))
print(np.sin(arr))
\end{lstlisting}
\end{examplebox}

\begin{warning}{* vs @ / 逐元素乘法 vs 矩阵乘法}
    \begin{itemize}[itemsep=0.3em]
        \item \texttt{arr * arr}：Element-wise multiplication / 逐元素乘法.
        \item \texttt{arr @ arr}：Matrix product (inner product for 1D vectors) / 矩阵相乘（一维时为内积）.
    \end{itemize}
    混淆这两个操作会导致完全不同的结果！
\end{warning}

\subsection{Aggregation Functions / 聚合函数与轴}

\begin{concept}{Reducing Dimensionality / 降维操作}
    Aggregation functions (min, max, sum, average) reduce dimensionality. They provide an \textbf{axis} argument to specify which dimension to reduce. \\
    聚合函数降低数组维度。\texttt{axis} 参数指定沿着哪个轴聚合。
\end{concept}

\begin{examplebox}{Aggregation with axis / 聚合与轴参数}
\begin{lstlisting}
two_dim_arr = np.random.randint(0, high=20, size=(3, 4))
print(two_dim_arr)              # 3x4 random array

# Aggregation over whole array / 对整个数组聚合
print(np.min(two_dim_arr))      # Minimum of ALL elements
print(np.max(two_dim_arr))      # Maximum of ALL elements
print(np.sum(two_dim_arr))      # Sum of ALL elements

# axis=0: Operation applied across rows (column-wise)
# 纵向聚合：按列计算，结果 shape 为 (4,)
print(np.min(two_dim_arr, axis=0))  # Min of each column

# axis=1: Operation applied across columns (row-wise)
# 横向聚合：按行计算，结果 shape 为 (3,)
print(np.average(two_dim_arr, axis=1))  # Avg of each row
\end{lstlisting}
\end{examplebox}

\begin{warning}{axis Meaning / axis 参数含义}
    \begin{itemize}
        \item \textbf{\texttt{axis=0}}：Operation across rows — reduces to one value per \textbf{column}. / 纵向（按列计算）.
        \item \textbf{\texttt{axis=1}}：Operation across columns — reduces to one value per \textbf{row}. / 横向（按行计算）.
    \end{itemize}
    Think of it as: the axis you specify is the one that \textbf{disappears}!
\end{warning}

\subsection{Combining Arrays / 数组合并}

\begin{examplebox}{concatenate, append, stack / 数组合并三剑客}
\begin{lstlisting}
a = np.arange(10)                   # [0 1 2 3 4 5 6 7 8 9]
b = np.arange(10)[::-1]             # [9 8 7 6 5 4 3 2 1 0]

# concatenate: joins a sequence of arrays
print(np.concatenate((a, b)))
# [0 1 2 3 4 5 6 7 8 9 9 8 7 6 5 4 3 2 1 0]

# append: uses concatenate internally
print(np.append(a, b))

# stack: stacks along a NEW axis (creates new dimension)
print(np.stack((np.arange(10), np.arange(10))))
# [[0 1 2 3 4 5 6 7 8 9]
#  [0 1 2 3 4 5 6 7 8 9]]
\end{lstlisting}
\end{examplebox}

\begin{warning}{Concatenation Creates Copies / 合并操作会创建副本}
    Operations like \texttt{np.append}, \texttt{np.concatenate}, and \texttt{np.stack} always require the whole array to be copied. It often makes more sense to allocate an array of the size you need \textbf{upfront} and then fill it.
\end{warning}

\subsection{Masking / 掩码过滤}

\begin{concept}{Boolean Indexing / 布尔索引}
    Logical arrays (arrays containing boolean values) can be used to index other arrays. These logical arrays are called \textbf{masks}. This is especially useful to filter data based on conditions. \\
    布尔数组可以用作索引来过滤数据，称为"掩码"。
\end{concept}

\begin{examplebox}{Masking Examples / 掩码过滤示例}
\begin{lstlisting}
arr = np.arange(1, 6)               # [1 2 3 4 5]

# Create a boolean mask
mask = np.array([True, False, True, False, True])
print(arr[mask])                    # [1 3 5] (only where mask is True)

# Generate mask from condition
print(arr < 5)                      # [ True  True  True  True False]
print(arr[arr < 3])                 # [1 2] (filtered by condition)
\end{lstlisting}
\end{examplebox}

\subsection{Advanced Indexing / 高级索引}

\begin{examplebox}{2D Array Indexing / 二维数组索引}
    The syntax works as \texttt{[row, column]}, or more precisely \texttt{[row\_start:row\_end, col\_start:col\_end]}.
\begin{lstlisting}
large_two_dim_arr = np.arange(81).reshape((9, 9))

# Select entire column / 选择整列
print(large_two_dim_arr[:, 2])      # All rows, column 2

# Select specific row and column range
print(large_two_dim_arr[1, 2:7])    # Row 1, columns 2 through 6

# Standard slicing with step / 带步长的切片
print(large_two_dim_arr[:, 2:7:2])  # All rows, columns 2, 4, 6
\end{lstlisting}
\end{examplebox}

\end{document}
